\chapter{Jaundice}

\section*{Definition}
Yellowish discoloration of the skin, sclerae, and mucous membranes due to elevated serum bilirubin levels, clinically detectable when bilirubin exceeds 2-3 mg/dL (34-51 µmol/L).

\section*{What Happens in Your Body}
Bilirubin accumulation results from pre-hepatic (hemolysis, ineffective erythropoiesis), hepatic (impaired hepatocyte uptake, conjugation, or excretion), or post-hepatic (biliary obstruction preventing bilirubin excretion into the gut) causes. Unconjugated (indirect) vs conjugated (direct) hyperbilirubinemia distinguished by laboratory testing guides diagnostic workup.

\section*{Causes}
\begin{itemize}
  \item Viral hepatitis (A, B, C, E)
  \item Alcoholic liver disease or cirrhosis
  \item Gallstones or bile duct obstruction
  \item Pancreatic cancer (head of pancreas)
  \item Hemolytic anemias (sickle cell, spherocytosis)
  \item Drug-induced liver injury (acetaminophen, statins, isoniazid)
  \item Primary biliary cholangitis
  \item Gilbert syndrome (benign unconjugated hyperbilirubinemia)
  \item Acute fatty liver of pregnancy
  \item Autoimmune hepatitis
\end{itemize}

\section*{When to See a Doctor}
See your doctor if:
\begin{itemize}
  \item Symptoms are severe or getting worse over time
  \item Any jaundice requires medical evaluation to identify underlying cause quickly
  \item You have other concerning symptoms such as fever, weight loss, or bleeding
\end{itemize}

\section*{Related Symptoms}
\begin{itemize}
  \item Dark urine
  \item Pale stools
  \item Abdominal pain
  \item Nausea
  \item Itching (pruritus)
\end{itemize}
\textbf{Important:} If this symptom persists, your doctor will check for serious underlying causes.

\section*{Self-Care / Home Management}
\begin{itemize}
  \item Seek medical evaluation promptly — jaundice always requires a medical assessment to determine the underlying cause; do not attempt self-treatment.
  \item Avoid alcohol completely until the cause of jaundice is identified, especially if liver injury or biliary obstruction is suspected.
  \item Check all medications and supplements with your doctor, as many drugs (including acetaminophen, statins, and herbal remedies) can cause or worsen liver injury.
  \item Rest and maintain adequate fluid intake, but note that dietary changes alone will not resolve jaundice without addressing the root cause.
\end{itemize}

\section*{How Your Doctor Will Evaluate You}
\begin{itemize}
  \item History focusing on onset (acute vs gradual), medication and alcohol use, travel, family history, and associated symptoms (pain, fever, pruritus, dark urine, pale stools).
  \item Liver function tests with fractionated bilirubin (total, direct/unconjugated), ALT, AST, ALP, GGT, and prothrombin time to classify the type and severity of jaundice.
  \item Abdominal ultrasound is the first-line imaging study to assess bile duct dilation, gallstones, liver parenchyma, and pancreatic pathology.
  \item Referral to a gastroenterologist or hepatologist is indicated for suspected chronic liver disease, biliary obstruction, or when the aetiology remains unclear after initial workup.
\end{itemize}

\section*{Treatment Options}
\begin{itemize}
  \item Obstructive jaundice (gallstones, stricture, tumour) requires endoscopic retrograde cholangiopancreatography (ERCP) or surgical intervention to relieve biliary obstruction.
  \item Viral hepatitis is managed with supportive care, antiviral therapy (hepatitis B or C), and avoidance of hepatotoxins including alcohol.
  \item Autoimmune hepatitis and primary biliary cholangitis are treated with immunosuppression (corticosteroids, azathioprine) or ursodeoxycholic acid respectively.
  \item Haemolytic jaundice is treated by addressing the underlying haemolytic disorder (e.g., splenectomy for hereditary spherocytosis, stopping offending drugs).
  \item Medications, lifestyle modifications, and physical therapies may be prescribed as appropriate.
  \item Follow-up ensures treatment effectiveness and allows timely adjustments to the care plan.
\end{itemize}

\section*{References}

\begin{thebibliography}{99}
\bibitem{jaundice:harrison-jaundice} Longo DL, et al. \emph{Harrison's Principles of Internal Medicine}. 21st ed. McGraw-Hill; 2022. Chapter 341: Jaundice.
\bibitem{jaundice:uptodate-jaundice} Roche SP, et al. Approach to the patient with jaundice. In: UpToDate, Post TW (Ed), UpToDate, Waltham, MA; 2025.
\bibitem{jaundice:beckingham} Beckingham IJ, et al. Jaundice: a clinical review. \emph{BMJ}. 2023;386:e074941.
\bibitem{jaundice:melichar} Melichar B, et al. Evaluation of the jaundiced patient. \emph{Am J Gastroenterol}. 2024;119(3):445-457.
\bibitem{jaundice:feldman} Feldman M, et al. \emph{Sleisenger and Fordtran's Gastrointestinal and Liver Disease}. 11th ed. Elsevier; 2023. Chapter 20: Jaundice.
\bibitem{jaundice:european-liver} European Association for the Study of the Liver. EASL guideline on the management of cholestatic liver diseases. \emph{J Hepatol}. 2024;80(2):303-326.
\end{thebibliography}
